% ---
% id: protocol-compiler/req:binary64-weights
% kind: req
% title: "Declared weights as binary64 input"
% summary: >-
%   A declared weight in either form that is not exactly a finite binary64
%   number is rejected as malformed input before it enters any sum or product.
%   This is the declared input type, not a failed predicate, and it has no
%   refusal string of its own.
% prominence: standard
% status: implemented
% uses: [def:weighted]
% assumes: []
% refs: []
% provenance: newly-derived
% implements: certiqs_key.protocol_compiler.weights:weighted_family
% lean:
% ---

The conditions of Definition~\ref{def:weighted} are decided on declared
numbers, so the declared weights have to be numbers of the type they are
decided on.

\begin{requirement}\label{req:binary64-weights}
The weights of either form of Definition~\ref{def:weighted} --- each $q_t$ of
the \tref{setting-independent form} and each entry of $P$ of the
\tref{setting-dependent form} --- are \tref{binary64 numbers}. The engine
rejects as malformed input a declared weight, in either form, that is not
exactly representable as a finite binary64 number, before that weight is used
in an accumulation or a product. This is the declared input type: the
rejection is not a failed summation predicate, and it adds no row of the gate
report and no refusal string.
\end{requirement}
